\documentclass[10pt,a4paper]{article}

% ==============================
% Packages
% ==============================
\usepackage[utf8]{inputenc}
\usepackage[T1]{fontenc}
\usepackage{amsmath,amssymb,amsfonts}
\usepackage{mathtools}
\usepackage{bm}
\usepackage{geometry}
\geometry{margin=1in}
\usepackage{hyperref}
\hypersetup{colorlinks=true, linkcolor=blue, citecolor=blue, urlcolor=blue}
\usepackage{natbib}
\usepackage{booktabs}
\usepackage{enumitem}
\usepackage{xcolor}

% ==============================
% Custom commands (adapt to your style)
% ==============================
\newcommand{\Msun}{M_{\odot}}
\newcommand{\OmHI}{\Omega_{\mathrm{HI}}}
\newcommand{\bHI}{b_{\mathrm{HI}}}
\newcommand{\Tbar}{\bar{T}_b}
\newcommand{\MHI}{M_{\mathrm{HI}}}
\newcommand{\cHI}{c_{\mathrm{HI}}}
\newcommand{\Tsys}{T_{\mathrm{sys}}}
\newcommand{\Tinst}{T_{\mathrm{inst}}}
\newcommand{\Tsky}{T_{\mathrm{sky}}}
\newcommand{\vc}{v_{\mathrm{c}}}
\newcommand{\vczero}{v_{\mathrm{c},0}}
\newcommand{\fHc}{f_{H,c}}
\newcommand{\rhocrit}{\rho_{\mathrm{crit}}}
\newcommand{\rhoHI}{\rho_{\mathrm{HI}}}
\newcommand{\utilde}{\tilde{u}}
\newcommand{\hmpc}{h^{-1}\,\mathrm{Mpc}}

% ==============================
\title{Cross-Correlating \texorpdfstring{H\textsc{i}}{HI} 21-cm and Fermi-LAT 
Gamma-Ray Maps:\\
A Comprehensive Guide to the Halo-Model Framework \\
and MeerKAT Instrumentation}

\author{PhD Research Notes --- Thread~2}
\date{\today}

% ==============================
\begin{document}
% \maketitle

% \begin{abstract}
% These notes provide a comprehensive, PhD-level treatment of the neutral hydrogen
% (H\textsc{i}) intensity-mapping side of the cross-correlation framework developed
% by Pinetti, Camera, Fornengo \& Regis (2020, arXiv:2001.02649; JCAP 07, 044).  We cover the
% astrophysics of post-reionisation H\textsc{i}, the Padmanabhan, Refregier \&
% Amara (2017) halo model, MeerKAT instrumental specifics, foreground contamination,
% and the data-processing pipeline.  For each topic we explain the physics, identify
% the modelling choices made in Pinetti et al., assess how well those choices have
% held up against five years of subsequent data (2020--2026), and point to key
% references.
% \end{abstract}

% \tableofcontents
% \newpage


% =====================================================================
\section{Where neutral hydrogen survives after reionisation}
\label{sec:post-reion-HI}
% =====================================================================

After reionisation completes at $z \approx 6$, the intergalactic medium is
overwhelmingly ionised (neutral fraction $\lesssim 10^{-4}$).  Surviving neutral
hydrogen resides almost exclusively in \emph{self-shielded, overdense gas clouds}
--- damped Lyman-$\alpha$ systems (DLAs) with column densities
$N_{\mathrm{HI}} \ge 2 \times 10^{20}\;\mathrm{cm}^{-2}$ --- within the
interstellar medium of galaxies hosted by dark matter haloes
\citep{Wolfe2005}.


\subsection{The critical halo-mass window}

Pinetti et al.\ state that haloes of mass
$5 \times 10^{9}$--$10^{12}\;\Msun$ contain most of the cosmic H\textsc{i}.
This range emerges from two physical boundary effects:

\paragraph{Low-mass cutoff ($M \lesssim 10^{8}\;\Msun$).}
Haloes below $\sim\!10^{8}\;\Msun$ have virial (circular) velocities
$\vc \lesssim 25$--$30\;\mathrm{km\,s}^{-1}$.  The metagalactic UV
photoionisation background heats the IGM to $T \sim 10^{4}\;\mathrm{K}$,
raising the Jeans mass and preventing gas from accreting into these shallow
potential wells \citep{Rees1986,Efstathiou1992,Pontzen2008}.  Simulations
confirm that haloes with $\vc < 30\;\mathrm{km\,s}^{-1}$ do not efficiently
host neutral gas \citep{Bagla2010}.

\paragraph{High-mass suppression ($M \gtrsim 10^{12}\;\Msun$).}
Haloes above $\sim\!10^{12}\;\Msun$ have virial temperatures exceeding
$10^{5}$--$10^{6}\;\mathrm{K}$, where infalling gas is shock-heated to the
hot, X-ray-emitting phase rather than cooling to form a neutral component.
AGN feedback in these massive systems further heats and expels gas.

The intervening ``sweet spot'' at $10^{10}$--$10^{12}\;\Msun$ contains haloes
deep enough to accrete gas despite the UV background, yet cool enough for that
gas to settle into a neutral phase.  ALFALFA abundance-matching analyses and
the conditional H\textsc{i} mass function both confirm this picture
\citep{Papastergis2013,Padmanabhan2017}.


\subsection{The 21-cm hyperfine transition}

The line arises from the spin-flip of the electron in the hydrogen $1S$
ground state between the triplet ($F=1$, parallel spins, statistical weight~3)
and singlet ($F=0$, antiparallel, weight~1) hyperfine levels, separated by
\begin{equation}
  \Delta E = 5.88 \times 10^{-6}\;\mathrm{eV}.
\end{equation}
The rest-frame frequency is
\begin{equation}
  \nu_{21} = 1420.405\;\mathrm{MHz}
  \qquad (\lambda = 21.106\;\mathrm{cm}),
\end{equation}
and the spontaneous emission coefficient is
$A_{10} = 2.85 \times 10^{-15}\;\mathrm{s}^{-1}$, corresponding to a natural
lifetime of $\sim\!10^{7}\;\mathrm{yr}$.

In the post-reionisation regime, the spin temperature $T_s$ greatly exceeds
$T_{\mathrm{CMB}}$ via collisional coupling and the Wouthuysen--Field
(Ly$\alpha$ pumping) effect, so the 21-cm signal is always seen in
\emph{emission} against the CMB.  The brightness temperature then becomes
independent of $T_s$.  The observed frequency maps directly to redshift via
\begin{equation}
  \label{eq:nu-z}
  \nu_{\mathrm{obs}} = \frac{1420.405}{1+z}\;\mathrm{MHz},
\end{equation}
which is the foundation of H\textsc{i} intensity mapping as a
three-dimensional probe of large-scale structure
\citep[Pinetti et al.\ Eq.~3.1]{Bharadwaj2004}.


\subsection{The H\textsc{i} density parameter \texorpdfstring{$\OmHI(z)$}{Omega\_HI(z)}}

Defined as
\begin{equation}
  \label{eq:OmHI}
  \OmHI(z) = \frac{\rhoHI(z)}{\rho_{\mathrm{crit},0}},
\end{equation}
this parameter sets the overall amplitude of the H\textsc{i} signal.  In the
halo model it is computed as (Pinetti et al.\ Eq.~3.2)
\begin{equation}
  \label{eq:OmHI-halomodel}
  \OmHI = \frac{1}{\rhocrit}
  \int \frac{\mathrm{d}n}{\mathrm{d}M}\,\MHI(M,z)\,\mathrm{d}M.
\end{equation}

Observational constraints come from multiple probes.  Local 21-cm emission
surveys such as HIPASS give $\OmHI h_{70} \approx 3.8 \times 10^{-4}$ at
$z \sim 0$ \citep{Zwaan2005}, while the ALFALFA $\alpha$.100 catalogue gives
$\OmHI h_{70} \approx (3.9\text{--}4.3) \times 10^{-4}$ \citep{Jones2018}.
DLA absorbers from the SDSS-III sample yield
$\OmHI \approx (0.8\text{--}1.0) \times 10^{-3}$ at $z \sim 2$--3.5
\citep{Noterdaeme2012}, and the GGG survey extends to $z \sim 4.9$ with
$\OmHI = (0.98 \pm 0.20) \times 10^{-3}$ \citep{Crighton2015}.  Intensity
mapping cross-correlations such as GBT $\times$ WiggleZ constrain
$\OmHI \bHI$ at $z \sim 0.8$ \citep{Switzer2013,Masui2013}.

A commonly used empirical fit is $\OmHI(z) \propto (1+z)^{0.6}$, indicating
a gradual increase from $\sim\!4 \times 10^{-4}$ at $z=0$ to
$\sim\!10^{-3}$ at $z \sim 2$--5.  The cross-correlation signal amplitude
$C_\ell^{\mathrm{HI}\times\gamma}$ scales \emph{linearly} with $\OmHI$
through $\Tbar$, so the $\sim\!20$--30\% observational uncertainty in $\OmHI$
propagates directly into a comparable fractional uncertainty in the predicted
signal.


\subsection{Mean brightness temperature \texorpdfstring{$\Tbar(z)$}{Tbar(z)}}

In the post-reionisation limit ($T_s \gg T_{\mathrm{CMB}}$), the mean 21-cm
brightness temperature is (Pinetti et al.\ Eq.~3.4)
\begin{equation}
  \label{eq:Tbar}
  \boxed{
  \Tbar(z) \approx 188\,h\;\OmHI(z)\,
  \frac{(1+z)^{2}}{E(z)}\;\mathrm{mK},
  }
\end{equation}
where $E(z) = H(z)/H_0$.  Typical values are $\sim\!55\;\mu\mathrm{K}$ at
$z = 0$, $\sim\!140\;\mu\mathrm{K}$ at $z = 1$, and
$\sim\!240\;\mu\mathrm{K}$ at $z = 2$, increasing with redshift because both
$\OmHI$ and $(1+z)^{2}/E(z)$ grow.  These sub-millikelvin amplitudes --- four
to five orders of magnitude below the astrophysical foregrounds --- are the
fundamental challenge of intensity mapping.


% =====================================================================
\section{The H\textsc{i}--halo mass relation and its parameter space}
\label{sec:MHI-M}
% =====================================================================

Pinetti et al.\ adopt the parametric relation of
\citet{Padmanabhan2017} (their Eq.~3.7):
\begin{equation}
  \label{eq:MHI}
  \boxed{
  \MHI(M, z) = \alpha\,\fHc\,M\,
  \left(\frac{M}{10^{11}\,h^{-1}\,\Msun}\right)^{\!\beta}
  \exp\!\left[-\left(\frac{\vczero}{\vc(M,z)}\right)^{\!3}\right],
  }
\end{equation}
where $\fHc = (1 - Y_p) \approx 0.76$ is the cosmic hydrogen mass fraction
and $\vc(M,z)$ is the halo circular velocity.


\subsection{Physical meaning of each parameter}

\paragraph{Normalisation $\alpha$.}
Controls the overall fraction of baryonic hydrogen in neutral form per halo.
It directly scales $\OmHI$ and hence $\Tbar$.

\paragraph{Power-law slope $\beta$.}
Negative ($\beta \approx -0.58$ to $-0.69$), encoding the decline of
H\textsc{i} efficiency with halo mass: more massive haloes convert a smaller
fraction of their baryons to neutral gas, reflecting virial shock-heating and
AGN feedback.

\paragraph{Minimum circular velocity $\vczero$.}
Sets the low-mass cutoff.  Haloes with $\vc < \vczero$ have their H\textsc{i}
content exponentially suppressed, encoding photoionisation heating that
prevents gas accretion in shallow potential wells.


\subsection{Parameter calibration}

\citet{Padmanabhan2017} fit these parameters via MCMC (using CosmoHammer)
simultaneously to: (i)~the ALFALFA H\textsc{i} mass function at $z \sim 0$
\citep{Martin2010}; (ii)~DLA column-density distributions at $z \sim 2$--5
\citep{Noterdaeme2012,Crighton2015}; (iii)~$\OmHI$ measurements across
$z \sim 0$--5; (iv)~DLA incidence rate $\mathrm{d}N/\mathrm{d}X$ and DLA bias
$b_{\mathrm{DLA}}$; (v)~small-scale clustering from ALFALFA
\citep{Martin2012,Papastergis2013}; and (vi)~intensity mapping constraints at
$z \sim 0.8$ \citep{Switzer2013}.

The specific values adopted by Pinetti et al.\ are
\begin{equation}
  \alpha = 0.176, \qquad
  \beta  = -0.69, \qquad
  \vczero = 101.61\;\mathrm{km\,s}^{-1}.
\end{equation}
These differ notably from the Padmanabhan et al.\ (2017) published best-fits
for the altered-NFW profile ($\alpha = 0.17$, $\beta = -0.55$,
$\vczero = 37.2\;\mathrm{km\,s}^{-1}$) or the exponential profile
($\alpha = 0.09$, $\beta = -0.58$, $\vczero = 36.3\;\mathrm{km\,s}^{-1}$).
The notably larger $\vczero = 101.61\;\mathrm{km\,s}^{-1}$ implies a more
aggressive low-mass cutoff, placing the H\textsc{i} minimum-mass halo at
roughly $10^{10}\;\Msun$ rather than $\sim\!10^{8}$--$10^{9}\;\Msun$.  This
choice shifts the effective $\bHI$ upward while reducing the total $\OmHI$
contribution from low-mass haloes, and warrants careful verification against
the original parameter tables.


\subsection{Alternative models}

The principal competitors are cosmological hydrodynamical simulations.
\textbf{IllustrisTNG} \citep{VillaescusaNavarro2018} assigns H\textsc{i} in
post-processing using the \citet{Rahmati2013} self-shielding prescription; it
agrees well with $\OmHI$ at $z \sim 0$ but shows more complex $\MHI(M)$ with
significant scatter and environmental dependence (assembly bias) not captured
by the parametric model.  \textbf{SIMBA} \citep{Dave2019,Dave2020} uses
meshless-finite-mass hydrodynamics with torque-limited BH accretion and bipolar
kinetic AGN jets, and tends to produce lower H\textsc{i} masses in
intermediate-mass haloes than IllustrisTNG due to different AGN feedback
prescriptions.  \textbf{MUFASA} \citep{Dave2017} is the predecessor of SIMBA
with simpler feedback.

All models agree reasonably at $z \sim 0$ but diverge by factors of
2--5 at $z > 1$, particularly in $\OmHI$ contributions from haloes above
$10^{12}\;\Msun$.  The CHIME auto-power spectrum interpretation
\citep{CHIME2026} found 3.1--4.0$\sigma$ disagreement with IllustrisTNG
predictions at $z \sim 1$, likely due to inadequate modelling of nonlinear
redshift-space clustering of H\textsc{i} in the simulations.


\subsection{Sensitivity analysis}

On large scales, $C_\ell^{\mathrm{HI}\times\gamma} \propto \Tbar \times \bHI
\times W_\gamma \times P_m$.  Since $\Tbar \propto \alpha$ linearly,
\emph{$\alpha$ is the parameter that most directly scales the overall signal
amplitude}.  The slope $\beta$ redistributes H\textsc{i} across halo masses,
affecting both $\OmHI$ and the effective $\bHI$ in a correlated but partially
cancelling manner.  The cutoff $\vczero$ has the most complex effect:
increasing it reduces $\OmHI$ but raises $\bHI$ (by removing low-bias,
low-mass haloes), so the net impact on the product $\OmHI \times \bHI$ depends
on the steepness of the halo mass function.

Fisher-forecast analyses \citep{Padmanabhan2019} show the H\textsc{i}
auto-power spectrum is most sensitive to $\vczero$ and $\beta$ as
astrophysical parameters, while the cross-correlation amplitude uncertainty
is dominated by the combined product $\OmHI \times \bHI$.


% =====================================================================
\section{The H\textsc{i} density profile within haloes}
\label{sec:profile}
% =====================================================================

Pinetti et al.\ model the H\textsc{i} radial density within haloes using the
altered NFW profile with a thermal core (their Eq.~3.9):
\begin{equation}
  \label{eq:rhoHI-profile}
  \boxed{
  \rhoHI(r) = \frac{\rho_0\,r_s^{3}}
  {\bigl(r + 0.75\,r_s\bigr)\bigl(r + r_s\bigr)^{2}},
  }
\end{equation}
where $r_s = R_{\mathrm{vir}}/\cHI$ is the scale radius.  This profile was
introduced by \citet{MallerBullock2004} for multiphase halo gas and
subsequently adopted for DLA studies by Barnes \& Haehnelt (2010, 2014) and
incorporated into the halo model by \citet{Padmanabhan2017}.


\subsection{Why a core, not a cusp}

Unlike collisionless dark matter, which forms the familiar $r^{-1}$ NFW cusp,
neutral hydrogen gas has \emph{thermal pressure support}.  The warm neutral
medium at $T \sim 10^{4}\;\mathrm{K}$ exerts gas pressure that resists
gravitational concentration toward the halo centre, producing a flat core at
$r \lesssim 0.75\,r_s$.  Additionally, gas retains angular momentum during
infall, settling into disc-like configurations rather than cuspy spheroids.
The profile transitions smoothly to the standard NFW $r^{-3}$ envelope at
large radii ($r \gg r_s$).  The $0.75\,r_s$ core scale arises naturally from
thermal equilibrium considerations in multiphase cooling simulations
\citep{Cole2000,MallerBullock2004}.


\subsection{H\textsc{i} concentration \texorpdfstring{$\cHI$}{c\_HI}}

\citet{Padmanabhan2017} parametrise the H\textsc{i} concentration as
\begin{equation}
  \cHI(M, z) = c_{\mathrm{HI},0}\,
  \left(\frac{M}{10^{11}\,\Msun}\right)^{-0.109}
  \times \frac{4}{(1+z)^{0.13}},
\end{equation}
with a best-fit $c_{\mathrm{HI},0} \approx 139$ and redshift exponent
$\gamma \approx 0.13$.  This is roughly
\emph{ten times larger} than typical dark matter concentrations
($c_{\mathrm{DM}} \sim 10$--20; \citealt{Maccio2007}).  The physical reason
is straightforward: baryonic gas can \emph{dissipate energy} through radiative
cooling (atomic line emission, Compton cooling), losing kinetic energy and
contracting to much smaller radii than collisionless dark matter.  For a
Milky-Way-like galaxy, the H\textsc{i} disc extends to $\sim\!15$--30~kpc
while the virial radius is $\sim\!200$--350~kpc --- a ratio directly
mirroring the high $\cHI$.


\subsection{Normalisation via mass conservation}

The central density $\rho_0$ is uniquely determined by requiring (Pinetti
et al.\ Eq.~3.10):
\begin{equation}
  \label{eq:mass-conservation}
  \int_{0}^{R_{\mathrm{vir}}} \rhoHI(r)\,4\pi r^{2}\,\mathrm{d}r
  = \MHI(M, z).
\end{equation}
This is simply mass conservation.


\subsection{Fourier transform and angular-scale dependence}

The normalised Fourier transform that enters the power spectrum calculation is
(Pinetti et al.\ Eq.~3.14)
\begin{equation}
  \label{eq:uHI}
  \utilde_{\mathrm{HI}}(k \mid M, z) = \frac{1}{\MHI}
  \int_{0}^{R_{\mathrm{vir}}} \rhoHI(r)\,
  \frac{\sin(kr)}{kr}\,4\pi r^{2}\,\mathrm{d}r.
\end{equation}
Because $\cHI \gg c_{\mathrm{DM}}$, the H\textsc{i} profile is far more
compact and $\utilde_{\mathrm{HI}} \approx 1$ out to much higher $k$ than the
DM NFW profile.  On \textbf{large scales} ($\ell \lesssim 100$), the 2-halo
term dominates and both profiles give $\utilde \to 1$, so the profile shape is
irrelevant.  The profile choice becomes significant for the \textbf{1-halo
term at $\ell \gtrsim 200$--500}, corresponding to
$k \gtrsim 0.1$--$0.3\;\hmpc$ at $z \sim 0.3$--1.

For the Pinetti et al.\ cross-correlation forecast, which spans
$\ell \sim 10$--1000, profile details matter primarily at the high-$\ell$ end.


\subsection{Robustness from simulations and observations}

IllustrisTNG shows that mean H\textsc{i} profiles are approximately universal
across mass and redshift but exhibit large halo-to-halo scatter
\citep{VillaescusaNavarro2018}.  Resolved H\textsc{i} observations from the
THINGS survey \citep{Walter2008} and WHISP show that azimuthally averaged
profiles flatten or decline near the centre and fall exponentially in the outer
regions --- better described by an \emph{exponential disc} than a modified NFW
at low redshift, as acknowledged by \citet{Padmanabhan2017}.  For the
cross-correlation analysis at $z \sim 0.1$--0.5, the modified NFW remains a
reasonable approximation since the profile primarily affects the 1-halo term
at high $\ell$ where instrumental noise is large.


% =====================================================================
\section{H\textsc{i} bias: how neutral hydrogen traces large-scale structure}
\label{sec:bias}
% =====================================================================

The effective H\textsc{i} bias in the halo model is (Pinetti et al.\ Eq.~3.6):
\begin{equation}
  \label{eq:bHI}
  \boxed{
  \bHI(z) = \frac{1}{\bar{\rho}_{\mathrm{HI}}}
  \int \mathrm{d}M\;\frac{\mathrm{d}n}{\mathrm{d}M}\;
  \MHI(M, z)\;b(M, z).
  }
\end{equation}
This is a \emph{mass-weighted} average of the halo bias, where the weighting
is by H\textsc{i} content rather than number or total mass.


\subsection{Predicted and measured values}

\begin{table}[htbp]
\centering
\caption{H\textsc{i} bias measurements and predictions at different redshifts.}
\label{tab:bHI}
\begin{tabular}{@{}lll@{}}
\toprule
Redshift & $\bHI$ (predicted) & Observational constraint \\
\midrule
$z \sim 0$   & 0.65--0.85 & ALFALFA: $\bHI \approx 0.8$ \citep{Martin2012} \\
$z \sim 0.4$ & 0.8--1.0   & MeerKAT $\times$ WiggleZ:
  $\OmHI \bHI r = (0.86 \pm 0.16) \times 10^{-3}$ \citep{Cunnington2023} \\
$z \sim 0.8$ & 0.9--1.1   & GBT $\times$ WiggleZ:
  $\OmHI \bHI r = (0.43 \pm 0.07) \times 10^{-3}$ \citep{Masui2013} \\
$z \sim 2.3$ & 1.5--2.2   & DLA bias:
  $b_{\mathrm{DLA}} = 2.17 \pm 0.2$ \citep{FontRibera2012} \\
\bottomrule
\end{tabular}
\end{table}

The H\textsc{i} bias \emph{increases monotonically with redshift}.  At
$z \sim 0$, H\textsc{i} traces relatively underdense regions (blue, late-type
galaxies), giving $\bHI < 1$.  By $z \sim 2$--3, H\textsc{i}-hosting haloes
become rarer peaks in the density field with higher bias, reaching
$\bHI \sim 2$--3.


\subsection{Scale dependence}

Linear (scale-independent) bias breaks down at
$k \sim 0.3\;\hmpc$ at $z \le 3$ \citep{VillaescusaNavarro2018,Castorina2017}.
Using the Limber approximation ($\ell \sim k \chi(z)$) for $z \sim 0.3$--1
where $\chi \sim 1$--3~Gpc, this corresponds to $\ell \sim 300$--900.  For the
Pinetti et al.\ multipole range of $\ell = 10$--1000, the regime
$\ell \lesssim 100$--300 is safely in the linear bias regime, while
$\ell > 300$--500 enters the nonlinear regime where the 1-halo term and
scale-dependent bias become important.


\subsection{Impact on the cross-correlation}

On large scales, $C_\ell^{\mathrm{HI}\times\gamma} \propto \bHI$ linearly,
so a 20\% uncertainty in $\bHI$ translates directly to a $\sim\!20\%$
uncertainty in the astrophysical cross-correlation amplitude.  Current
constraints on $\bHI$ have 20--50\% uncertainties depending on redshift, making
it one of the dominant astrophysical systematics alongside $\MHI(M, z)$ and
$\OmHI$.  However, joint measurement of H\textsc{i} auto- and cross-power
spectra can break the $\OmHI$--$\bHI$ degeneracy, since
\begin{equation}
  P_{\mathrm{auto}} \propto \bHI^{2}, \qquad
  P_{\mathrm{cross}} \propto \bHI.
\end{equation}


% =====================================================================
\section{Observational validation: from ALFALFA through MeerKAT and CHIME}
\label{sec:validation}
% =====================================================================

\subsection{The ALFALFA survey}

The Arecibo Legacy Fast ALFA survey used the 7-beam ALFA receiver on the 305-m
Arecibo telescope to perform a blind extragalactic H\textsc{i} survey over
$\sim\!7000\;\mathrm{deg}^{2}$ at $z < 0.06$.  The final $\alpha$.100
catalogue contains $\sim\!31{,}500$ extragalactic H\textsc{i} sources
\citep{Haynes2018}, spanning H\textsc{i} masses from $\sim\!10^{6}$ to
$\sim\!10^{10.8}\;\Msun$.  \citet{Martin2012} and \citet{Papastergis2013}
measured the correlation function and power spectrum of H\textsc{i}-selected
galaxies, finding that H\textsc{i} sources cluster more weakly than optically
selected samples, consistent with $\bHI \approx 0.8$ at $z \sim 0$.


\subsection{Pinetti et al.\ Figure~1 comparison}

The halo-model prediction agrees well on \textbf{large scales}
($k < 1\;\hmpc$), where the 2-halo term dominates, but systematically
\textbf{underestimates clustering on small scales} (below
$\sim\!1\;\hmpc$).  This discrepancy arises because the standard halo model
with a single continuous H\textsc{i} profile per halo does not capture
satellite galaxy contributions, gas stripping in group environments, or other
baryonic physics that redistribute H\textsc{i} on sub-halo scales.
\citet{Padmanabhan2023} showed that adding finger-of-God effects causes
$\gtrsim\!10\%$ deviations at $k \gtrsim 0.1\;h\;\mathrm{cMpc}^{-1}$, and
that satellite contributions enter primarily at $k > 10\;\hmpc$.


\subsection{The revolution of 2023--2026}

Since Pinetti et al.'s 2020 paper, the H\textsc{i} intensity mapping field has
undergone a transformation from upper limits to high-significance detections.

\paragraph{MeerKAT.}
The landmark result is the \textbf{7.7$\sigma$ cross-correlation} with WiggleZ
galaxies at $z \approx 0.40$--0.46 \citep{Cunnington2023}, measuring
$\OmHI \bHI r = (0.86 \pm 0.16) \times 10^{-3}$.  Paul et al.\ (2023)
achieved the first interferometric auto-power spectrum detection at $z = 0.32$
and 0.44, at 8.0$\sigma$ and 11.5$\sigma$ respectively on nonlinear scales
($0.3 < k < 8\;\mathrm{Mpc}^{-1}$).  \citet{Carucci2025} confirmed the
cross-correlation with improved multi-scale PCA cleaning, finding
$\OmHI \bHI r = (0.93 \pm 0.17) \times 10^{-3}$ at $\sim\!6\sigma$.  The
MeerKLASS deep field \citep{MeerKLASS2025} reached thermal noise of
$\sim\!1.21\;\mathrm{mK}$ --- the first time thermal noise is
\emph{subdominant to H\textsc{i} fluctuations} at
$k \lesssim 0.15\;\hmpc$.

\paragraph{CHIME.}
CHIME detected 21-cm emission via stacking on eBOSS catalogues at
$z = 0.78$--1.43 \citep{CHIME2023}, achieving 7.1$\sigma$ (LRG),
5.7$\sigma$ (ELG), and \textbf{11.1$\sigma$ (QSO)} detections.
Cross-correlation with the Ly$\alpha$ forest at $z \approx 2.3$ yielded a
\textbf{9$\sigma$} detection --- the highest-redshift 21-cm emission
measurement to date \citep{CHIME2024}.  The first \textbf{auto-power spectrum
detection at $z \sim 1$} with 12.5$\sigma$ significance was announced in late
2025 \citep{CHIME2025}, covering $k = 0.4$--$1.5\;\hmpc$ at $z = 1.01$--1.34.

\paragraph{FAST.}
FAST remains in the systematics mitigation phase for cosmological intensity
mapping, characterising $1/f$ noise (knee frequency
$\sim\!6 \times 10^{-4}\;\mathrm{Hz}$ after SVD subtraction;
\citealt{Hu2021}).

\paragraph{Does the Padmanabhan model hold up?}
The measured values of $\OmHI \times \bHI$ from MeerKAT and CHIME are
\emph{broadly consistent} with Padmanabhan et al.\ predictions on large
scales.  However, the CHIME $z \sim 1$ interpretation paper found
3.1--4.0$\sigma$ tension with IllustrisTNG \citep{CHIME2026}, suggesting that
simulations underpredict nonlinear H\textsc{i} clustering in redshift space.
The parametric halo model remains a viable framework, but extensions
incorporating improved redshift-space distortion modelling and satellite
contributions are needed for the precision era.


% =====================================================================
\section{MeerKAT as an intensity mapping instrument}
\label{sec:MeerKAT}
% =====================================================================

\subsection{Array configuration}

MeerKAT comprises \textbf{64 dishes} of \textbf{13.5~m diameter}
(offset Gregorian, unblocked aperture) in the Upper Karoo, South Africa.
About 70\% of dishes lie within a dense 1-km core (Gaussian dispersion
$\sim\!300$~m), with the remainder extending to maximum baselines of
$\sim\!8$~km.  Frequency coverage spans the UHF band (580--1015~MHz,
H\textsc{i} at $z = 0.4$--1.45) and L-band (900--1670~MHz, H\textsc{i} at
$z = 0.0$--0.58), each with 4096 frequency channels.  In practice, the L-band
cosmological window is severely limited by RFI to only 971--1023~MHz
($z \approx 0.39$--0.46), motivating the shift to UHF for MeerKLASS, which
provides $\sim\!40\times$ larger comoving volume \citep{Cunnington2025}.


\subsection{Single-dish vs.\ interferometric mode}

These modes access complementary angular scales.

\paragraph{Single-dish (auto-correlation) mode.}
Each dish operates as a scanning total-power radiometer, measuring large
angular scales ($\ell \lesssim$ few hundred) critical for cosmological scales.
The scanning strategy has dishes slewing at 5~arcmin/s in azimuth at constant
elevation, with noise-diode injection every 20~s for relative gain calibration.
The combined 64-dish system effectively multiplies integration time by 64.

\paragraph{Interferometric mode.}
2016 baselines from 64 dishes resolve structure on small scales but filter
out modes below $\ell_{\min}$ set by the shortest baseline ($\sim\!29$~m).


\subsection{System temperature and noise}

The noise model follows Pinetti et al.\ Eq.~3.18 for single-dish mode.  The
system temperature is
\begin{equation}
  \label{eq:Tsys}
  \Tsys = \Tinst + \Tsky, \qquad
  \Tsky = 60\left(\frac{300\;\mathrm{MHz}}{\nu}\right)^{2.55}\;\mathrm{K},
\end{equation}
where $\Tinst \sim 18$--30~K (cryogenic receivers; $\sim\!18$~K at L-band
centre per SARAO specifications, $\sim\!30$~K used in forecasts including
spillover) and $\Tsky$ is dominated by Galactic synchrotron emission.
At 1~GHz, $\Tsky \approx 3.3$~K gives $\Tsys \sim 20$--25~K; at 600~MHz,
$\Tsky \approx 10$~K gives $\Tsys \sim 25$--40~K.

The thermal noise power spectrum scales as
\begin{equation}
  P_N \propto \frac{\Tsys^{2}\,\Omega_{\mathrm{survey}}}
  {N_{\mathrm{dish}}\,t_{\mathrm{obs}}\,\delta\nu\,N_{\mathrm{pol}}}.
\end{equation}
Achieved noise levels are 1.2--1.4$\times$ the theoretical limit:
\citet{Wang2021} reported $\sim\!2\;\mathrm{mK}$ in the pilot survey, while
the MeerKLASS deep field (2025) reached $\sim\!1.21\;\mathrm{mK}$.


\subsection{Beam and angular resolution}

The single-dish beam is approximately Gaussian with FWHM (Pinetti et al.\
Eq.~3.17):
\begin{equation}
  \theta_{\mathrm{FWHM}} \approx \frac{\lambda}{D_{\mathrm{dish}}}.
\end{equation}
At 1~GHz, $\theta_{\mathrm{FWHM}} \approx 1.27^{\circ}$; at 700~MHz,
$\approx 1.82^{\circ}$; at 580~MHz, $\approx 2.19^{\circ}$.
Frequency-dependent beam structure (chromaticity) couples with foregrounds and
can bias signal reconstruction \citep{Spinelli2021}; the \texttt{eidos}
package provides holographic beam models for more careful treatment.


\subsection{Interferometric noise}

The interferometric noise (Pinetti et al.\ Eq.~3.19) depends on the baseline
density $n(u)$.  At 1~GHz:
\begin{equation}
  \ell_{\min} \approx 2\pi \times \frac{29\;\mathrm{m}}{0.3\;\mathrm{m}}
  \approx 607, \qquad
  \ell_{\max} \approx 2\pi \times \frac{8000}{0.3} \approx 167{,}000.
\end{equation}
The dense core provides excellent short-baseline sensitivity, while the sparse
outer baselines probe very small angular scales.  The crossover multipole where
interferometric noise beats single-dish noise occurs at
$\ell \sim 200$--600, depending on frequency.  For MeerKAT, cosmological
scales ($k \lesssim 0.2\;\hmpc$) are accessible only in single-dish mode; the
interferometer covers the nonlinear regime.


\subsection{MeerKLASS status update (2026)}

The MeerKAT Large Area Synoptic Survey \citep{Santos2017} originally proposed
$\sim\!4000\;\mathrm{deg}^{2}$ over $\sim\!4000$~hours.  As of early 2026,
MeerKLASS has accumulated $> 650$ hours of UHF-band data reaching nearly
$3000\;\mathrm{deg}^{2}$, making it the largest spectroscopic survey of
large-scale structure in the Southern hemisphere.  Under the Extended Large
Project (XLP) allocation, the survey targets
$10{,}000\;\mathrm{deg}^{2}$ by 2028 with $\sim\!2500$ total hours.  The
\textbf{MeerKAT+ upgrade} adds 13--16 SKA-design dishes (first antenna handed
over February 2024), extending the maximum baseline to $\sim\!17$~km and
increasing raw sensitivity by $\sim\!50\%$.  MeerKAT+ will eventually
integrate into \textbf{SKA1-MID} (197 dishes total).


% =====================================================================
\section{Foregrounds: enormous but spectrally distinguishable}
\label{sec:foregrounds}
% =====================================================================

\subsection{The scale of the problem}

Astrophysical foregrounds at MeerKAT frequencies are \textbf{4--5 orders of
magnitude brighter} than the H\textsc{i} cosmological signal.  Galactic
synchrotron emission dominates, with brightness temperatures of $\sim\!1$--6~K
at high Galactic latitudes at 1~GHz (tens to hundreds of K near the plane),
following a spectral index $\beta \approx -2.5$ to $-2.9$.  Galactic free-free
emission contributes $\sim\!1$--10~K with a flatter spectral index
$\beta \approx -2.1$.  Extragalactic point sources (radio galaxies, AGN) add
both clustered and Poisson components.  The total foreground brightness is of
order \textbf{Kelvin}, versus an H\textsc{i} signal of order
\textbf{0.1--0.3~mK rms}.


\subsection{Why removal is possible}

All dominant foregrounds have \emph{smooth power-law spectra} that vary slowly
with frequency, while the cosmological 21-cm signal has rich structure along
the frequency (line-of-sight) direction.  Foregrounds therefore occupy only a
few modes in the frequency dimension, making them separable in principle.  In
interferometric data, the chromatic response of baselines confines foreground
power to a \textbf{wedge-shaped region} in $(k_\perp, k_\parallel)$ Fourier
space, leaving a ``clean window'' at high $k_\parallel$ for signal extraction.

The primary limitations are: (i)~\textbf{frequency-dependent beam}
(chromaticity) that mixes foreground power into signal modes;
(ii)~\textbf{calibration errors} that introduce spectral structure mimicking
H\textsc{i}; and (iii)~\textbf{polarisation leakage} combined with Galactic
Faraday rotation that creates non-smooth spectral features.


\subsection{Removal techniques}

The MeerKAT community employs several approaches.

\paragraph{PCA/SVD (blind).}
Decomposes the data matrix (frequency $\times$ pixel) into eigenmodes,
removing the first $N_{\mathrm{fg}} = 4$--30 modes identified as foregrounds.
\citet{Cunnington2023} used $N_{\mathrm{fg}} = 30$.  \textbf{Multi-scale PCA}
\citep{Carucci2025} applies PCA with different $N_{\mathrm{fg}}$ at different
angular scales, optimising foreground removal while preserving cosmological
signal.

\paragraph{Gaussian Process Regression (GPR).}
Models data as smooth foreground (long frequency correlation length) plus
signal+noise (short correlation length).  Better at recovering radial
($k_\parallel$) power spectrum than PCA \citep{Soares2022}.

\paragraph{SDecGMCA.}
Sparse component separation that jointly performs beam deconvolution and source
separation \citep{Gkogkou2025}, achieving $\lesssim\!5\%$ accuracy at
$\ell = 10$--200 under realistic beam conditions --- outperforming standard
methods.

\paragraph{Foreground avoidance.}
Used in interferometric mode (delay spectrum technique; \citealt{Paul2023}),
measuring the power spectrum only in the foreground-free region of Fourier
space.

All methods cause some \textbf{signal loss}, quantified via the transfer
function
\begin{equation}
  T(k) = \frac{P_{\mathrm{recovered}}(k)}{P_{\mathrm{injected}}(k)},
\end{equation}
constructed by injecting mock H\textsc{i} signal into real data, re-running the
cleaning pipeline, and measuring the recovery fraction.  Despite $\sim\!50\%$
signal loss on the largest scales, the power spectrum can be reconstructed to
sub-percent accuracy using this transfer function approach.


\subsection{The cross-correlation advantage for
\texorpdfstring{H\textsc{i} $\times$ gamma rays}{HI x gamma rays}}

This is perhaps the most important point for the Pinetti et al.\ analysis.
Foregrounds in H\textsc{i} maps (Galactic synchrotron, free-free) and
foregrounds in Fermi-LAT gamma-ray maps (Galactic diffuse gamma-ray emission
from cosmic-ray interactions) have \textbf{completely different physical
origins} and are therefore largely \textbf{uncorrelated}.  Foreground residuals
in one map do not systematically bias the cross-correlation measurement ---
they increase variance but not the mean of the cross-power estimate.

This has been demonstrated observationally: \citet{Cunnington2019} showed that
foregrounds have little impact on optical spectroscopic cross-correlations, with
reconstructed cross-power within 8.5\% even at scales below beam resolution.
The practical success of cross-correlation detections --- MeerKAT $\times$
WiggleZ at 7.7$\sigma$, CHIME $\times$ eBOSS at 11.1$\sigma$, CHIME $\times$
Ly$\alpha$ at 9$\sigma$ despite foreground residuals 6--10$\times$ thermal
noise --- constitutes strong empirical evidence.


\subsection{Spurious cross-correlation risks}

The primary concern is \textbf{radio-loud AGN that are also gamma-ray
emitters}.  The Fermi-LAT 4FGL-DR4 catalogue contains 7194 sources, of which
$>60\%$ are blazars --- and essentially all blazars are radio-loud by
definition.  Mitigation requires: (i)~masking all 4FGL source positions and
their error ellipses from both maps; (ii)~masking known bright radio sources
from continuum catalogues (NVSS, SUMSS, VLASS) above a flux threshold; and
(iii)~statistically estimating the residual contribution from unresolved
sources below detection thresholds using source-count models.


% =====================================================================
\section{From raw MeerKAT data to science-ready intensity maps}
\label{sec:pipeline}
% =====================================================================

\subsection{Calibration}

Standard interferometric calibration achieves bandpass stability within
$\sim\!3\%$ over 3~hours.  Intensity mapping requires stability at the
$\sim\!10^{-5}$ level to separate the $\mu$K H\textsc{i} signal from
Kelvin-scale foregrounds.  The MeerKLASS single-dish pipeline
\citep{Wang2021} models total-power data as a sum of components: calibrator
temperature, diffuse foreground, elevation-dependent terrestrial emission,
noise-diode signal, and receiver residual.  Calibration uses periodic
noise-diode injection (1.8~s every 20~s) for relative gain tracking, with
absolute calibration from celestial sources (3C~273, Pictor~A) every
$\sim\!1.5$ hours.

\paragraph{$1/f$ noise.}
Correlated temporal gain fluctuations are the critical systematic for
single-dish mode.  \citet{Li2021} measured a knee frequency of
$\sim\!3 \times 10^{-3}\;\mathrm{Hz}$ after 2-SVD-mode subtraction, implying
adequate gain stability over timescales of a few hundred seconds.  The $1/f$
noise is highly correlated in frequency and hence removable via SVD/PCA, but
aggressive cleaning risks signal loss.

\paragraph{Polarisation leakage.}
Stokes $Q, V \to I$ leakage combined with Faraday rotation introduces
non-smooth spectral structure requiring careful correction.


\subsection{RFI excision}

Despite the Karoo's radio-quiet zone designation, RFI from
telecommunications satellites (especially in L-band $\sim\!1150$--1300~MHz),
distant mobile towers, and aircraft transponders remains significant.  The
\citet{Wang2021} pipeline uses a 3-round flagging approach: (i)~strong RFI via
the \texttt{SEEK} package (SumThreshold algorithm), flagging 40--60\% across
the full L-band but only \textbf{5--25\% in target cosmological windows};
(ii)~per-channel outlier removal in time-ordered data; and (iii)~map-based
identification of residual weak RFI.  The UHF band suffers less severe
contamination, with RFI restricted to narrow frequency chunks.  Additional
tools include Tricolour (GPU-accelerated flagger) and deep learning approaches
(R-Net architecture).


\subsection{Map-making}

Calibrated time-ordered data (Stokes~$I$ in Kelvin) are projected into sky
maps using optimal map-making:
\begin{equation}
  \hat{m} = \bigl(\mathbf{A}^{\!\top}\mathbf{N}^{-1}\mathbf{A}\bigr)^{-1}
  \mathbf{A}^{\!\top}\mathbf{N}^{-1}\mathbf{d},
\end{equation}
where $\mathbf{A}$ is the pointing matrix and $\mathbf{N}$ the noise
covariance.  The MeerKLASS pilot used $0.3^{\circ}$ square pixels in flat-sky
approximation; recent work employs Cartesian regridding of spherical maps to
avoid wide-sky biases.  HEALPix pixelisation is used in simulations and some
analysis frameworks.  Foreground subtraction (PCA/GPR/GMCA) then operates on
these frequency--pixel data cubes.  The transfer function is computed at this
stage via mock signal injection.


\subsection{Software ecosystem}

The MeerKAT community has developed a rich set of publicly available tools.
\textbf{CARACal} (formerly MeerKATHI) is an end-to-end containerised
calibration and imaging pipeline for interferometric data, using the Stimela
framework \citep{Jozsa2020}.  \textbf{katdal} is a Python data access library
for MeerKAT Visibility Format (MVF) files, handling HDF5/RDB chunk stores and
conversion to CASA MeasurementSets.  \textbf{katcali} is the
MeerKLASS-specific single-dish calibration pipeline for auto-correlation
intensity mapping, handling RFI flagging, noise-diode calibration, absolute
calibration, and map-making.  \textbf{oxkat} provides semi-automated
calibration and imaging scripts covering the full workflow from flagging through
self-calibration \citep{Heywood2020}.  \textbf{katbeam} and \textbf{eidos}
provide primary beam models (simplified cosine-taper and holographic,
respectively).  \textbf{MeerFish} is a Fisher forecast code released alongside
\citet{Cunnington2025}.


\subsection{Public data products}

MeerKLASS has released the \textbf{first public H\textsc{i} intensity mapping
data cube} from the 2019 pilot survey via
\url{https://meerklass.org/data.html}, covering
$153^{\circ} < \mathrm{RA} < 172^{\circ}$,
$-1^{\circ} < \mathrm{Dec} < 8^{\circ}$ as calibrated intensity cubes after
map-making but before foreground removal.  UHF and L-band continuum data
releases provide complementary radio source catalogues.  MIGHTEE deep-field
data in COSMOS, E-CDFS, ELAIS-S1, and XMM-LSS fields support interferometric
intensity mapping analyses.


% =====================================================================
\section{Conclusion: a framework that has aged well but needs updating}
\label{sec:conclusion}
% =====================================================================

The Pinetti et al.\ (2020) framework rests on a solid theoretical foundation
that has been broadly validated by the subsequent explosion of observational
results.  Three key findings emerge from this assessment.

First, the Padmanabhan et al.\ (2017) H\textsc{i} halo model remains the
most practical analytic framework for cross-correlation predictions, but it
requires extensions for the precision era --- particularly improved modelling
of nonlinear redshift-space distortions, satellite contributions, and the
3.1--4.0$\sigma$ tension with IllustrisTNG identified by CHIME at $z \sim 1$.

Second, the cross-correlation advantage for foreground mitigation has been
\emph{spectacularly confirmed} by multiple experiments: MeerKAT, CHIME, and
GBT cross-correlations all succeed even when auto-correlations remain
foreground-dominated.  This bodes well for the H\textsc{i} $\times$ gamma-ray
cross-correlation, where the physical disconnect between radio and gamma-ray
foregrounds is even stronger than between radio and optical foregrounds.

Third, the dominant uncertainty in the predicted signal amplitude comes from
the product $\OmHI \times \bHI$, constrained to $\sim\!15$--20\% at
$z \sim 0.4$ by MeerKAT but still poorly known at $z > 1$.  For this
analysis, the specific parameter values adopted by Pinetti et al.\
(particularly $\vczero = 101.61\;\mathrm{km\,s}^{-1}$) should be carefully
checked against the original Padmanabhan et al.\ tables and tested for
consistency with the latest MeerKAT and CHIME constraints.

The MeerKLASS UHF-band survey, now exceeding $3000\;\mathrm{deg}^{2}$,
combined with 14~years of Fermi-LAT data (4FGL-DR4), places the first actual
measurement of the H\textsc{i} $\times$ gamma-ray cross-correlation within
realistic reach.


% =====================================================================
% BIBLIOGRAPHY
% =====================================================================
% Replace the thebibliography below with \bibliography{yourfile} and a
% .bib file when integrating into your document.  The entries below are
% provided for self-contained compilation.

\begin{thebibliography}{99}

\bibitem[Bagla et al.(2010)]{Bagla2010}
Bagla, J.\,S., Khandai, N., \& Datta, K.\,K.\ 2010, MNRAS, 407, 567

\bibitem[Barnes \& Haehnelt(2014)]{Barnes2014}
Barnes, L.\,A.\ \& Haehnelt, M.\,G.\ 2014, MNRAS, 440, 2313

\bibitem[Barnes \& Haehnelt(2010)]{Barnes2010}
Barnes, L.\,A.\ \& Haehnelt, M.\,G.\ 2010, MNRAS, 403, 870

\bibitem[Bharadwaj \& Ali(2004)]{Bharadwaj2004}
Bharadwaj, S.\ \& Ali, S.\,S.\ 2004, MNRAS, 352, 142

\bibitem[Carucci et al.(2025)]{Carucci2025}
Carucci, I.\,P.\ et al.\ 2025, A\&A, in press

\bibitem[Castorina \& Villaescusa-Navarro(2017)]{Castorina2017}
Castorina, E.\ \& Villaescusa-Navarro, F.\ 2017, MNRAS, 471, 1788

\bibitem[CHIME Collaboration(2023)]{CHIME2023}
CHIME Collaboration\ 2023, ApJ, 947, 16

\bibitem[CHIME Collaboration(2024)]{CHIME2024}
CHIME Collaboration\ 2024, ApJ, 963, 23

\bibitem[CHIME Collaboration(2025)]{CHIME2025}
CHIME Collaboration\ 2025, arXiv:2511.19620

\bibitem[CHIME Collaboration(2026)]{CHIME2026}
CHIME Collaboration\ 2026, arXiv:2603.25680

\bibitem[Cole et al.(2000)]{Cole2000}
Cole, S.\ et al.\ 2000, MNRAS, 319, 168

\bibitem[Crighton et al.(2015)]{Crighton2015}
Crighton, N.\,H.\,M.\ et al.\ 2015, MNRAS, 452, 217

\bibitem[Cunnington et al.(2019)]{Cunnington2019}
Cunnington, S.\ et al.\ 2019, MNRAS, 488, 5452

\bibitem[Cunnington et al.(2023)]{Cunnington2023}
Cunnington, S.\ et al.\ 2023, MNRAS, 518, 6262

\bibitem[Cunnington et al.(2025)]{Cunnington2025}
Cunnington, S.\ et al.\ 2025, MNRAS, submitted

\bibitem[Dav{\'e} et al.(2017)]{Dave2017}
Dav{\'e}, R.\ et al.\ 2017, MNRAS, 467, 115

\bibitem[Dav{\'e} et al.(2019)]{Dave2019}
Dav{\'e}, R.\ et al.\ 2019, MNRAS, 486, 2827

\bibitem[Dav{\'e} et al.(2020)]{Dave2020}
Dav{\'e}, R.\ et al.\ 2020, MNRAS, 497, 146

\bibitem[Efstathiou(1992)]{Efstathiou1992}
Efstathiou, G.\ 1992, MNRAS, 256, 43P

\bibitem[Font-Ribera et al.(2012)]{FontRibera2012}
Font-Ribera, A.\ et al.\ 2012, JCAP, 11, 059

\bibitem[Gkogkou et al.(2025)]{Gkogkou2025}
Gkogkou, A.\ et al.\ 2025, A\&A, in press

\bibitem[Haynes et al.(2018)]{Haynes2018}
Haynes, M.\,P.\ et al.\ 2018, ApJ, 861, 49

\bibitem[Heywood(2020)]{Heywood2020}
Heywood, I.\ 2020, ascl:2009.003

\bibitem[Hu et al.(2021)]{Hu2021}
Hu, W.\ et al.\ 2021, arXiv:2109.06447

\bibitem[Jones et al.(2018)]{Jones2018}
Jones, M.\,G.\ et al.\ 2018, MNRAS, 477, 2

\bibitem[J{\'o}zsa et al.(2020)]{Jozsa2020}
J{\'o}zsa, G.\,I.\,G.\ et al.\ 2020, ASPC, 527, 635

\bibitem[Li et al.(2021)]{Li2021}
Li, Y.\,C.\ et al.\ 2021, MNRAS, submitted

\bibitem[Macci{\`o} et al.(2007)]{Maccio2007}
Macci{\`o}, A.\,V.\ et al.\ 2007, MNRAS, 378, 55

\bibitem[Maller \& Bullock(2004)]{MallerBullock2004}
Maller, A.\,H.\ \& Bullock, J.\,S.\ 2004, MNRAS, 355, 694

\bibitem[Martin et al.(2010)]{Martin2010}
Martin, A.\,M.\ et al.\ 2010, ApJ, 723, 1359

\bibitem[Martin et al.(2012)]{Martin2012}
Martin, A.\,M.\ et al.\ 2012, ApJ, 750, 38

\bibitem[Masui et al.(2013)]{Masui2013}
Masui, K.\,W.\ et al.\ 2013, ApJ, 763, L20

\bibitem[MeerKLASS Collaboration(2025)]{MeerKLASS2025}
MeerKLASS Collaboration\ 2025, MNRAS, 537, 3632

\bibitem[Noterdaeme et al.(2012)]{Noterdaeme2012}
Noterdaeme, P.\ et al.\ 2012, A\&A, 547, L1

\bibitem[Padmanabhan et al.(2017)]{Padmanabhan2017}
Padmanabhan, H., Refregier, A., \& Amara, A.\ 2017, MNRAS, 469, 2323

\bibitem[Padmanabhan et al.(2019)]{Padmanabhan2019}
Padmanabhan, H.\ et al.\ 2019, MNRAS, 485, 4060

\bibitem[Padmanabhan et al.(2023)]{Padmanabhan2023}
Padmanabhan, H.\ et al.\ 2023, MNRAS, submitted

\bibitem[Papastergis et al.(2013)]{Papastergis2013}
Papastergis, E.\ et al.\ 2013, ApJ, 776, 43

\bibitem[Paul et al.(2023)]{Paul2023}
Paul, S.\ et al.\ 2023, MNRAS, 525, 3439

\bibitem[Pinetti et al.(2020)]{Pinetti2020}
Pinetti, E., Camera, S., Fornengo, N., \& Regis, M.\ 2020, JCAP, 07, 044

\bibitem[Pontzen et al.(2008)]{Pontzen2008}
Pontzen, A.\ et al.\ 2008, MNRAS, 390, 1349

\bibitem[Rahmati et al.(2013)]{Rahmati2013}
Rahmati, A.\ et al.\ 2013, MNRAS, 430, 2427

\bibitem[Rees(1986)]{Rees1986}
Rees, M.\,J.\ 1986, MNRAS, 218, 25P

\bibitem[Santos et al.(2017)]{Santos2017}
Santos, M.\,G.\ et al.\ 2017, arXiv:1709.06099

\bibitem[Soares et al.(2022)]{Soares2022}
Soares, P.\,S.\ et al.\ 2022, MNRAS, 510, 5872

\bibitem[Spinelli et al.(2021)]{Spinelli2021}
Spinelli, M.\ et al.\ 2021, MNRAS, submitted

\bibitem[Switzer et al.(2013)]{Switzer2013}
Switzer, E.\,R.\ et al.\ 2013, MNRAS, 434, L46

\bibitem[Villaescusa-Navarro et al.(2018)]{VillaescusaNavarro2018}
Villaescusa-Navarro, F.\ et al.\ 2018, ApJ, 866, 135

\bibitem[Walter et al.(2008)]{Walter2008}
Walter, F.\ et al.\ 2008, AJ, 136, 2563

\bibitem[Wang et al.(2021)]{Wang2021}
Wang, J.\ et al.\ 2021, MNRAS, 505, 3698

\bibitem[Wolfe et al.(2005)]{Wolfe2005}
Wolfe, A.\,M., Gawiser, E., \& Prochaska, J.\,X.\ 2005, ARA\&A, 43, 861

\bibitem[Zwaan et al.(2005)]{Zwaan2005}
Zwaan, M.\,A.\ et al.\ 2005, MNRAS, 359, L30

\end{thebibliography}

\end{document}
